% Copyright (c) 2026 Twin Vector Labs LLC.
% All rights reserved.
\documentclass[11pt]{article}

\usepackage[margin=1in]{geometry}
\usepackage{amsmath,amssymb,amsthm}
\usepackage{mathtools}
\usepackage{enumitem}
\usepackage[colorlinks=true,linkcolor=blue,urlcolor=blue]{hyperref}

\newtheorem{principle}{Principle}
\theoremstyle{definition}
\newtheorem{remark}{Remark}

\newcommand{\C}{\mathbb{C}}
\newcommand{\Z}{\mathbb{Z}}
\newcommand{\rep}[1]{\mathbf{#1}}
\newcommand{\triv}{\mathbf{1}}
\newcommand{\Sg}{\Sigma}
\newcommand{\lsq}{\ell^2}
\DeclareMathOperator{\tr}{tr}

\title{\bf Design Choices for the Tripartite Proton Junction $W_{ABC}$\\
\large distinct holes, the relaxation fix, and why the singlet needs complex geometry}
\author{Tessera --- cobordism subsystem (\#396)}
\date{}

\begin{document}
\maketitle

\begin{abstract}
The companion note showed that a \emph{bipartite} merge cannot produce the
color-singlet proton: the singlet is the three-index alternating invariant
$\varepsilon_{ijk}$, absent from $\rep 3\otimes\rep 3$. We therefore build the
proton on a \emph{trivalent} junction $W_{ABC}$ with boundary
$\partial W = \psi_A\sqcup\psi_B\sqcup\psi_C$. This note records the four design
choices made in that construction and shows that each is forced by the same fact
--- the singlet is irreducibly a \emph{three-body} object. (1) The three inputs
must enter as \emph{distinct} register holes (independent $1$-cycles), never as
stacked layers of one shared register, or the relaxation averages them. (2) The
metric relaxation must \emph{gate} the state-residual gradient once the inputs are
carried, or an ill-conditioned step explodes. (3) A van Raamsdonk metric built
from \emph{pairwise} mutual information cannot select the singlet, because the
singlet's pairwise MI is maximal and \emph{uniform} --- its content is a
three-party correlation. (4) The $\Z_3$ phases of $[1,\omega,\omega^2]$ live in
the complex boundary state; to let the \emph{geometry} carry them the metric must
go \emph{complex}, i.e.\ Lorentzian --- which simultaneously realizes the photon
channel (null worldline edges).
\end{abstract}

\section{The junction}

The register is the homology of a holed surface. For the proton we take one
connected base surface --- a $2$-frequency geodesic icosahedron ($S^2$, $42$
vertices) --- minus $12$ vertex-disjoint hole triangles grouped into \textbf{four
windows of three}: $A,B,C$ (the three color inputs) and $R$ (the emergent
result). Extruding $\times I$ (\texttt{prismCells}) gives a connected $3$-complex
$W$ with $b_1 = 11$, a genuine valid manifold (every triangle in $\le 2$
tetrahedra, \texttt{dualComplexValid}) --- not a weld. Each window's three color
holes carry a color amplitude triple; the carried object is a color rep, and
charge is the oriented sum $\Sg(\psi)=\sum_i\psi_i$ ($\Sg=0$ is color-neutral).

\section{Choice 1 --- distinct holes, not stacked layers}

A single \emph{shared} register (the bipartite construction, and the linear
multi-layer staircase) fails for a precise reason.

\begin{principle}[Shared registers average]
In a product complex $M\times I$ the copy of a given hole-circle on each layer is
the \emph{same} homology class (the layers cobound the tube wall). Pinning three
inputs on one shared class asks a single period to equal three distinct targets;
\texttt{residualForPeriods} returns the least-squares solution --- the
\emph{average} --- which retains the symmetric (sextet) component and is never the
totally antisymmetric singlet.
\end{principle}

This is why the ``$b_1$ question'' is a false binary: both the $b_1=2$ linear
staircase $(S^2-3\text{ holes})\times[0,3]$ and the $b_1=3$ four-holed sphere
stack the inputs on one shared cycle, and both average. What carries the singlet
is the opposite: the three inputs as \emph{independent} $1$-cycles --- distinct
holes on the base surface. Confinement is then \emph{conservation at the
junction}: on a connected surface-minus-holes the global Stokes relation forces
\[
  \sum_{\text{all holes}} \big(\text{induced period}\big) \;=\; 0
  \qquad\Longrightarrow\qquad
  \Sg_R \;=\; -\,\Sg_{\text{inputs}} .
\]
Three color-neutral inputs ($\Sg=0$ each) therefore force a color-neutral result.
Measured on the built junction: the three neutral pairs carry \emph{exactly}
($r_U\sim 10^{-27}$, independent cycles --- no over-determination); the emergent
result is neutral for neutral inputs ($|\Sg_R|\sim 0$) and carries the net charge
for colored inputs ($|\Sg_R|\approx 3$); and the construction is $S_3$
(color-relabel) invariant.

\section{Choice 2 --- gate the state-residual gradient}

The junction relaxation initially stalled at iteration $0$ (no improving step)
where the smaller bipartite relaxed fine. The cause is a direct consequence of
Choice~1.

The relaxation minimizes $\;\beta\,\lVert\nabla_I S\rVert^2 + r_{\text{state}}\;$
over the interior edge lengths by a Gauss--Newton / Levenberg--Marquardt step.
Because the inputs are carried \emph{exactly} on independent cycles,
$r_{\text{state}}\sim 10^{-27}$ on \emph{any} metric --- so its analytic gradient
is numerical noise. Folded into the step through the (ill-conditioned) action
Gauss--Newton Hessian, that noise gradient explodes the step,
$\lVert\delta\rVert\sim 10^{13}$, and every line-search trial diverges.

\begin{remark}[The fix]
Fold the state-residual gradient into the step \emph{only} while
$r_{\text{state}}$ is above its convergence floor; once the state term is
converged, minimize the action alone (and guard the line search against a
degenerate-step throw). The bipartite/operator paths
($r_{\text{state}}\gg\varepsilon$) are unchanged; the junction now relaxes
($0\to 40^+$ iterations, $\lVert\nabla S\rVert^2:\,315\to 14$). The action-only
step is provably a descent direction --- the explosion was entirely the spurious
state gradient.
\end{remark}

\section{Choice 3 --- why a pairwise-MI van Raamsdonk metric cannot select the singlet}

The principled metric is van Raamsdonk's: geometry from entanglement,
\[
  \lsq_{uv} \;=\; \big(-\log(I_{uv}/I_{\max})\big)^2 ,
\]
with $I_{uv}$ the mutual information of the regions joined by edge $(u,v)$. The
hope was that seeding the metric from the color-singlet's entanglement would
refine the result toward the pure singlet. It cannot, for a sharp reason.

\begin{principle}[The singlet's pairwise MI is uniform]
For \emph{both} candidate three-qutrit singlets --- the totally antisymmetric
$\varepsilon$ state $\sum_{\sigma}\operatorname{sgn}(\sigma)\,|\sigma(RGB)\rangle$
and the $\Z_3$ clock state $|RRR\rangle+\omega|GGG\rangle+\omega^2|BBB\rangle$ ---
every single-party marginal is maximally mixed and every two-party reduced state
has spectrum $\{0,\tfrac13,\tfrac13,\tfrac13\}$. Hence
\[
  S(\rho_X)=\log 3,\qquad
  I(X{:}Y)=\log 3\quad\text{for all parties }X,Y,
\]
i.e.\ the pairwise mutual information is \emph{maximal and uniform}.
\end{principle}

A metric whose every (non-bulk) edge gets the same $I=\log 3$ is uniform: it has
no structure to distinguish the singlet, and indeed it does not move the result
(singlet-overlap $0.737$ with the VR metric vs.\ $0.744$ with random jitter). The
singlet's distinguishing content --- the relative $\omega$-phases --- is a
genuine \emph{three-party} correlation, invisible to any pairwise marginal. This
is the same irreducible-three-body fact that forced the tripartite junction in the
first place. Worse, the van Raamsdonk length is manifestly \emph{real}
($\lsq\ge 0$), so a VR metric cannot encode a complex phase under any
circumstances.

\section{Choice 4 --- the Lorentzian lever (and the photon)}

The proton singlet $[1,\omega,\omega^2]$, $\omega=e^{2\pi i/3}$, is intrinsically
\emph{complex}. A real metric yields real harmonic $1$-cochains, whose carried
transport is real and cannot represent the $\omega$-phases; in our construction
those phases enter only through the complex boundary inputs (the $\Z_3$-phased
neutral pairs already reach $\approx 0.74$ overlap). To let the \emph{geometry}
carry the phases --- and lift that ceiling --- the metric itself must be complex.

The model already has exactly one source of complex geometry: \emph{Lorentzian}
signature. A forward-time worldline edge (endpoints on different layers) carries
$\lsq<0$ (timelike), and the dual Regge action then goes complex,
$\operatorname{Im}S\ne 0$, so its harmonic spectrum --- and the carried transport
--- becomes complex. The same mechanism realizes the photon hypothesis: an edge
driven through $\lsq\to 0$ is \emph{null} (lightlike), a photon worldline. Thus
\emph{the lever that lets the geometry carry the singlet phases is the same lever
that emits the photon}.

A first probe (Lorentzian seed, before relaxation) is consistent with this:
making the cross-layer worldlines timelike nudges the overlap up
($0.659\to 0.671$) and the result toward neutrality ($|\Sg_R|:\,0.064\to 0.027$).
The full effect requires relaxing the complex (Lorentzian) geometry, where the
action's imaginary part constrains the scale degree of freedom the real part
leaves loose --- the work now in progress.

\section{Summary}

\begin{center}
\begin{tabular}{@{}p{0.30\linewidth}p{0.62\linewidth}@{}}
\hline
\textbf{Choice} & \textbf{Forced because the singlet is three-body}\\
\hline
Distinct holes (3 independent cycles) & a shared register averages three inputs into the symmetric part; independent cycles + Stokes give $\Sg_R=-\Sg_{\text{inputs}}$ (confinement by conservation).\\[2pt]
Gate the state gradient & independent cycles carry the inputs \emph{exactly}, so $r_{\text{state}}$ is floored and its gradient is noise that explodes the step; drop it once converged.\\[2pt]
Reject pairwise-MI VR & the singlet's pairwise MI is uniform ($\log 3$); its content is a 3-party phase correlation a real, pairwise metric cannot see or carry.\\[2pt]
Lorentzian (complex) geometry & the $\omega$-phases need a complex metric; timelike worldlines make $S$ complex, and the same null-edge limit is the photon.\\
\hline
\end{tabular}
\end{center}

\noindent
The junction is a valid, charge-conserving, $S_3$-respecting manifold whose
relaxation is fixed and which transports the complex boundary to a result
$\approx 74\%$ aligned with the ideal singlet. The remaining gap is a
\emph{phase} gap, and the route to closing it is complex (Lorentzian) geometry,
not a finer real metric.

\end{document}
